\documentclass[10pt,a4paper]{amsart}
\usepackage[margin=19mm]{geometry}
\usepackage{amsmath,amssymb,mathtools}
\usepackage[scr=rsfs,cal=euler]{mathalfa}
\usepackage{xfrac}
\usepackage[unicode,hidelinks,bookmarks=false]{hyperref}
\newcommand{\ie}{i.e.\ }
\newcommand{\cf}{cf.\ }
\newcommand{\resp}{respectively\ }
\newcommand{\Subsection}{\subsection}
\providecommand{\nobreakdash}{\nobreak\hbox{-}}
\setlength{\emergencystretch}{2em}
\multlinegap0pt
\usepackage{amssymb}
\usepackage{bm}
\usepackage{tikz}
\usepackage{cancel}
\newtheorem{theorem}{Theorem}
\title{A Hybridizable Neural Time Integrator for Stable Autoregressive Forecasting}
\author{Brooks Kinch, Xiaozhe Hu, Yilong Huang, Martine Dyring Hansen, Sunniva Meltzer, Nathaniel Donald Hamlin, David Sirajuddin, Eric C. Cyr,, Nathaniel Trask}
\date{Source: arXiv:2604.21101v2 (CC BY 4.0)}
\begin{document}
\maketitle

\noindent\emph{Source and licence.} A Hybridizable Neural Time Integrator for Stable Autoregressive Forecasting. Version arXiv:2604.21101v2 (CC BY 4.0), distributed by arXiv under the Creative Commons Attribution 4.0 International licence, http://creativecommons.org/licenses/by/4.0/, as recorded in the arXiv licence metadata for this exact version and shown on the abstract page https://arxiv.org/abs/2604.21101v2. Full licence notice: \texttt{licenses/arxiv-2604.21101-CC-BY-4.0.txt}.\par\smallskip

\noindent\textit{Editorial scope.} Counted unit: the article's theorem thm:hodgetheory (source0.tex lines 341--355). The article's complete original proof of it is delivered with it (source0.tex lines 357--374, one \texttt{proof} environment of 18 lines). No mathematics is written, repaired or re-derived: the statement and the proof are the article's own lines, in the article's own notation, and no retained line is substituted or re-flowed.  No further source-local context is needed: every cross-reference of every retained line resolves inside the two delivered spans.  Nothing was omitted from any retained span, so the package carries no omission record.  No retained line contains a citation, so the wrapper carries no bibliography and no bibliography entry.  No retained line contains a figure environment, an image-inclusion command or drawing code, so no image is delivered and the package declares no asset dependency.  Editorial formatting only, in addition: an AMS \texttt{amsart} preamble that restates the article's own declarations and macros used by the retained lines, namely the environments \texttt{theorem} on separate counters, together with the standard packages those lines need.  The retained environments print in this document as Theorem 1 in order of appearance, whereas the article prints the same items with the numbering of its own full document.  The complete native source of the article as distributed in the arXiv e-print for this exact version is bundled unchanged as \texttt{sources/198998/source0.tex}.
\begin{theorem}\label{thm:hodgetheory}
    Let $u,q \in \mathcal{Q}_h^i$, $J,v \in \mathcal{V}_h^i$, and $\lambda,\mu \in \mathcal{M}_h^i$. Then the following two identities, respectively, encode \textit{strong} and \textit{weak} time derivatives as products of the time derivative mass matrix and adjacency matrix.
    $$({J}',q) = \hat{q}^\intercal  \delta \hat{J}.$$
    $$(u,{v}') = \hat{v}^\intercal \delta^\intercal \hat{u}.$$
    Consider the following mixed Galerkin discretization of the Poisson problem with Dirichlet data. Find $(u,J) \in \mathcal{Q}_h^i\times \mathcal{V}_h^i$ such that for any $(q,v) \in \mathcal{Q}_h^i\times \mathcal{V}_h^i,$
    $$ (J,v) + (u,v') = \langle u_D,v \rangle, $$
    $$ (J',q) = (f,q),$$
    where $u_D \in \mathcal{M}_h^i$ is Dirichlet data and $f \in \mathcal{Q}_h^i$ is the forcing term. The standard Galerkin procedure yields the following system
    $$\mathbf{L} \hat{u} :=   \delta^\intercal \mathbf{M}_{\mathcal{V}}^{-1}\delta\hat{u} = \mathbf{M}_{\mathcal{Q}}\hat{f} + \mathbf{M}_\mathcal{V}^{-1}\mathbf{M}_\mathcal{M} \hat{u}_D.
 $$
    The stiffness matrix $\mathbf{L}$ corresponds to a discretized Hodge Laplacian operator, which is non-singular and symmetric positive definite with Poincar\'e inequality 
    $$\hat{u}^\intercal \mathbf{M}_{\mathcal{Q}} \hat{u} \leq C_p \hat{u}^\intercal \mathbf{L} \hat{u}, 
    $$ 
    where $C_p$ is independent of $h$.
\end{theorem}

\begin{proof}
Note that the time derivative operator $\frac{\mathrm{d}}{\mathrm{d}t} : \mathcal{V}_h^i \mapsto \mathcal{Q}_h^i$.  Then, its matrix representation $\mathbf{D} \in \mathbb{R}^{M \times (M+1)}$ is defined as 
\begin{align*}
\mathbf{D}_{ij} = (\frac{\mathrm{d}}{\mathrm{d}t}\psi_j^1, \phi_i^0) = 
\begin{cases}
-\frac{1}{h}, \quad  \text{if} \ j = i,  \\
\frac{1}{h}, \quad  \text{if} \ j=i+1, \\
0, \quad \text{otherwise.}
\end{cases}
\end{align*}
On the other hand, it is easy to compute that $\mathbf{M}_{\mathcal{Q}} = \operatorname{diag}(h, \cdots, h)$. Therefore, 
\begin{align*}
	(J', q) = \hat{q}^{\intercal} \mathbf{M}_{\mathcal{Q}} \mathbf{D} \hat{J} = \hat{q}^{\intercal} \delta \hat{J},
\end{align*} 
where we use the fact that $\mathbf{M}_{\mathcal{Q}} \mathbf{D} = \delta$ based on the calculation above.  Similarly, we have $(u, v') = \hat{v}^{\intercal}\delta^{\intercal} \hat{u}$ and the definition of $\mathbf{L}$ follows naturally.  

The discrete version of the Poincar\'e inequality follows directly from the fact that our choice of $\mathcal{V}_h^i \times \mathcal{Q}_h^i$ is an inf-sup stable pair for the mixed formulation in 1D. 
\end{proof}

\end{document}
